\section{Introduction}
\label{sec:introduction}

Modern communication networks face a wide spectrum of cyberattacks ranging from high-volume denial-of-service floods and automated reconnaissance scans to low-footprint shellcode injections, backdoors, and self-propagating worms. A network intrusion detection system (NIDS) inspects flow-level packet, protocol, and content attributes to distinguish benign traffic from specific attack categories~\cite{Moustafa2015}. Framing multi-category intrusion detection as a supervised classification problem introduces several practical challenges. Network traffic combines high-cardinality categorical protocol descriptors with skewed continuous rate and timing counters. Flow statistics also exhibit strong pairwise collinearity and severe class imbalance where benign flows and common attacks outnumber rare exploits by orders of magnitude.

Decision trees provide a natural baseline for intrusion detection because tree induction accommodates mixed categorical and numerical predictors without requiring feature rescaling or one-hot encoding, and represents predictions as explicit rule paths~\cite{Kelleher2020,Quinlan1986}. However, unrestricted decision trees grown greedily until leaf homogeneity exhibit high variance and memorise noisy or conflicting flow records. Two primary paradigms exist to address overfitting in tree-based classifiers. The first paradigm applies structural regularisation to a single tree through maximum depth constraints or post-pruning on held-out observations~\cite{Kelleher2020,Quinlan1993}. The second paradigm constructs an ensemble that combines multiple base learners through bootstrap aggregation, known as bagging, or sequential adaptive reweighting, known as boosting~\cite{Breiman1996,Freund1997,Zhu2009}.

The primary objective of the present study was to compare the predictive performance of a standalone decision tree with ensembles of decision trees for multiclass network intrusion detection on the University of New South Wales network-based 2015 (UNSW-NB15) dataset~\cite{Moustafa2015}. Rather than comparing an ensemble against a dissimilar model family such as nearest-neighbour or linear classifiers, the experimental design held the base learner family fixed across all experiments. Every standalone and ensemble model used the same custom iterative dichotomiser three (ID3) decision-tree implementation extended for continuous thresholds and reduced-error post-pruning. Holding the base learner constant ensured that observed differences in classification performance, class-level recall, and computational cost reflected the regularisation or aggregation mechanism rather than differences in feature representation.

To investigate the comparative objective systematically, the 257,673-row dataset was partitioned into fixed, stratified growth, pruning, validation, and reserved generalisation sets containing 123,683, 41,225, 41,234, and 51,531 rows, respectively. Four specific research questions guided the empirical progression:
\begin{enumerate}
    \item How did predictive performance scale as the training sample size increased, and how strongly did natural class scarcity in small samples constrain minority-attack recognition?
    \item How did maximum tree depth limits, reduced-error post-pruning, and class-balanced resampling interact to trade off overall accuracy against minority-class recall and precision?
    \item How did bagging and stagewise additive modelling using a multi-class exponential loss function (SAMME) behave as ensemble size, base-learner depth, and learning rate varied?
    \item Under a matched full-growth-pool budget and a shared hybrid class-balancing policy, how did an unpruned tree, a reduced-error-pruned tree, a bagging ensemble, and a SAMME boosting ensemble compare on the untouched generalisation set?
\end{enumerate}

The empirical results showed that ensemble aggregation did not automatically outperform a single regularised decision tree when macro-averaged F1-score across all ten traffic categories served as the primary metric. In learning-curve experiments from 2,000 to 50,000 rows, increasing the volume of original training data yielded much larger macro-averaged F1-score gains than ensembling at small sample sizes, while unweighted accuracy-based pruning on naturally skewed samples reduced macro-averaged F1-score by collapsing sparse minority branches. In the final comparison using all 123,683 growth rows under a hybrid class-balancing policy and pruning against all 41,225 dedicated pruning rows, the effect of pruning reversed. Training-set balancing allowed minority-attack subtrees to grow during induction, while reduced-error pruning on the natural pruning distribution removed overfitted leaves with high false-positive rates. On the reserved generalisation set, the pruned single tree achieved the highest mean macro-averaged F1-score of $0.5993$, the highest balanced accuracy of $63.02\%$, and the highest overall accuracy of $80.73\%$, outperforming 25-member bagging at $0.5747$, the unpruned tree at $0.5545$, and 25-member depth-eight SAMME boosting at $0.5408$. A supplementary validation extension confirmed that expanding bagging from 25 to 50 members produced no additional gain in validation macro-averaged F1-score.

The remainder of the report is organised into five sections. Section~\ref{sec:background} reviews the theoretical foundations of information-based decision trees, pruning strategies, and ensemble learning. Section~\ref{sec:methodology} details the dataset preprocessing, the custom decision-tree and ensemble implementations, algorithmic optimisations, and the hybrid class-balancing formulation. Section~\ref{sec:empirics} describes the four-way partitioning protocol, evaluation metrics, and the sequential design of the experimental stages. Section~\ref{sec:results} presents and analyses the empirical results across the development and generalisation sets. Finally, Section~\ref{sec:conclusion} summarises the conclusions, methodological limitations, and directions for future work.
